% ------------------------------------------------------------------------------
% Metodologia
% ------------------------------------------------------------------------------
\chapter{Metodologia}
\label{chap_metodologia}
Este capítulo descreve os procedimentos adotados para simular caminhadas aleatórias unidimensionais e comparar os resultados obtidos com a distribuição teórica do deslocamento final.

\section{Materiais Computacionais e Organização dos Dados}

As simulações foram realizadas em Python, utilizando \texttt{numpy} para geração de amostras e operações vetorizadas, \texttt{math} para o cálculo combinatório e \texttt{matplotlib} para visualização gráfica.

O experimento foi implementado em notebook Jupyter, disponível em \cite{github_random_walk}. O notebook está organizado para permitir reprodutibilidade completa dos resultados.

\section{Configuração do Experimento}

Foram definidos os seguintes parâmetros principais:
\begin{itemize}
    \item número de caminhadas: $10$;
    \item número de passos por caminhada: $100$;
    \item posição inicial: $S_0=0$;
    \item probabilidade de passo à direita: $1/2$;
    \item probabilidade de passo à esquerda: $1/2$.
\end{itemize}

Cada caminhada é construída a partir de uma sequência de 100 variáveis em $\{-1,+1\}$, com distribuição equiprovável. A posição em cada instante é obtida por soma cumulativa dos passos.

\section{Procedimento Experimental}
\label{sec_procedimento_experimental}

O procedimento computacional foi dividido em duas partes:
\begin{enumerate}
    \item \textbf{Simulação principal (10 caminhadas):} geração das 10 trajetórias com 100 passos cada, armazenamento das posições ao longo do tempo e obtenção dos deslocamentos finais.
    \item \textbf{Validação estatística (Monte Carlo):} geração de um conjunto ampliado de caminhadas (20000 simulações) para estimar probabilidades empíricas de deslocamento e compará-las com $P_n(d)$ teórico.
\end{enumerate}

Para a comparação empírico-teórica, foi calculada a frequência relativa de cada deslocamento final observado e, nos mesmos pontos, a probabilidade teórica via coeficiente binomial.

\section{Métricas e Produtos Gerados}

Foram produzidos os seguintes resultados:
\begin{itemize}
    \item gráfico das 10 trajetórias em função do número de passos;
    \item lista dos deslocamentos finais das 10 caminhadas;
    \item gráfico comparativo entre distribuição empírica (Monte Carlo) e distribuição teórica;
    \item erro absoluto máximo entre probabilidades empíricas e teóricas, como medida sintética de aderência.
\end{itemize}

A metodologia adotada permite verificar, de forma visual e quantitativa, a consistência entre simulação e teoria para o random walk simples em 1D.